\documentclass[10pt,a4paper]{amsart}
\usepackage[margin=19mm]{geometry}
\usepackage{amsmath,amssymb,mathtools}
\usepackage[scr=rsfs,cal=euler]{mathalfa}
\usepackage{xfrac}
\usepackage[unicode,hidelinks,bookmarks=false]{hyperref}
\newcommand{\ie}{i.e.\ }
\newcommand{\cf}{cf.\ }
\newcommand{\resp}{respectively\ }
\newcommand{\Subsection}{\subsection}
\providecommand{\nobreakdash}{\nobreak\hbox{-}}
\setlength{\emergencystretch}{2em}
\multlinegap0pt
\usepackage{bm}
\usepackage{bbm}
\usepackage{dsfont}
\usepackage{xspace}
\usepackage{cancel}
\usepackage{mathtools}
\usepackage{cleveref}
\usepackage[shortlabels]{enumitem}
\usepackage{amsthm}
\makeatletter
\@ifundefined{assumption}{\newtheorem{assumption}{Assumption}}{}
\@ifundefined{lemma}{\newtheorem{lemma}{Lemma}}{}
\@ifundefined{proposition}{\newtheorem{proposition}{Proposition}}{}
\@ifundefined{theorem}{\newtheorem{theorem}{Theorem}}{}
\makeatother
\DeclareMathOperator{\Var}{Var}
\def\Cov{{\textrm{Cov}}\,}
\def\R{\mathbb{R}}
\def\calN{\mathcal{N}}
\newcommand{\EE}[2][]{\mathbb{E}_{#1}\left[#2\right]}
\newcommand{\cv}{{\mathrm{CV}}}
\newcommand{\eqd}{\stackrel{d}{=}}
\newcommand{\norm}[1]{\left\| #1 \right\|}
\newcommand{\rd}{{\mathrm{d}}}
\newcommand{\tran}{^\intercal}
\title[On the optimality of antithetic randomization for cross-vali]{On the optimality of antithetic randomization for cross-validation}
\author{Srijan Chattopadhyay, Sifan Liu, Snigdha Panigrahi}
\date{Source: arXiv:2608.08089v1 (CC BY 4.0)}
\begin{document}
\maketitle

\noindent\emph{Source and licence.} On the optimality of antithetic randomization for cross-validation. Version arXiv:2608.08089v1 (CC BY 4.0), distributed under the Creative Commons Attribution 4.0 International licence, as recorded in the arXiv licence metadata for this exact version and shown on the abstract page https://arxiv.org/abs/2608.08089v1. Full rights notice: licenses/arxiv-2608.08089-CC-BY-4.0.txt.\par\smallskip

\noindent\textit{Editorial scope.} Counted unit: the Lemma whose source label is \texttt{lem:reducible:variance} in the article (source0.tex lines 942--974, source label \texttt{lem:reducible:variance}), delivered together with the article's own complete original proof of it (source0.tex lines 976--1160, one \texttt{proof} environment of 185 lines). No mathematics is written, repaired or re-derived: statement and proof are the article's own text line for line. Required context retained uncounted: assump:normal-randomization (lines 319--325); assump:weakdiff (lines 362--374); thm:alpha to 0 (lines 376--386); lem: recipe:AR (lines 403--412); lem:smooth-small-noise-limit (lines 646--683); equ:def:R\_alpha (lines 651--669). Every definition, display and label of those lines is retained unchanged. Omitted from this package and not redistributed: 1--318, 326--361, 375--375, 387--402, 413--645, 670--941, 975--975, 1161--2329. Nothing inside a retained excerpt is omitted or altered: every retained body is delivered exactly as the article's lines are written, so no omission receipt of any kind accompanies this package. Citations: the retained excerpts carry no citation command at all, so no bibliography entry of the article is transcribed and no reference is redistributed. Reference adaptation: none was needed and none was made; every reference inside the delivered unit resolves inside it, so no unresolved reference or citation is produced. Printed slips kept literal: no printed word, symbol, hypothesis or number of the counted statement or of its proof was altered. Layout and class: the article's own class is \texttt{article} with packages including graphicx, enumitem, xcolor, latexsym, amsmath, lscape, color, lineno, epstopdf, subfigure, hyperref, booktabs, multicol, multirow; the wrapper is the lane's pinned \texttt{amsart} editorial wrapper at 10pt on A4, which restates the theorem-like environments the retained text needs and adds the article's own macro definitions that the retained text needs (\texttt{\textbackslash Cov}, \texttt{\textbackslash EE}, \texttt{\textbackslash R}, \texttt{\textbackslash Var}, \texttt{\textbackslash calN}, \texttt{\textbackslash cv}, \texttt{\textbackslash eqd}, \texttt{\textbackslash norm}, \texttt{\textbackslash rd}, \texttt{\textbackslash tran}), with extra wrapper packages (bm, bbm, dsfont, xspace, cancel, mathtools, cleveref, enumitem). The delivered numbering is the unit's own. Assets: no retained span contains a figure environment, a graphics-inclusion command, a PDF-page inclusion, a table or a TikZ picture; no image file is copied into this package, the package declares no asset dependency and no image file is redistributed with it. Licence: version arXiv:2608.08089v1 of this article is distributed under the Creative Commons Attribution 4.0 International licence, as recorded in the arXiv licence metadata for this exact version and shown on the abstract page https://arxiv.org/abs/2608.08089v1. A full rights notice is delivered at licenses/arxiv-2608.08089-CC-BY-4.0.txt. Characters: every retained excerpt is pure ASCII, so the delivered engine, XeLaTeX with its default Latin Modern text font, reports no missing character.
\begin{assumption}[Exchangeable normal randomization]
\label{assump:normal-randomization}
  \begin{enumerate}
    \item[(i)]\label{assump:exch} \emph{Exchangeability:} \sloppy{$(\omega^{\pi(1)}, \ldots, \omega^{\pi(K)}) \eqd (\omega^1, \ldots, \omega^K)$} for every permutation $\pi$ of $[K]=\{1,\ldots,K\}$.
    \item[(ii)]\label{assump:marg} \emph{Marginal normality and equicorrelation:} $\omega^k\sim\calN(0,\sigma^2 I_n)$ for every $k\in[K]$, and $\Cov(\omega^k,\omega^{k'}) = \rho\sigma^2 I_n$ for every $k\neq k'$, where $\rho\in[-1/(K-1),1]$.
  \end{enumerate}
\end{assumption}

\begin{assumption}[Weak differentiability and moments at an inflated variance]
\label{assump:weakdiff}
    All components $g_i$ ($1 \le i \le n$) of $g$ are weakly differentiable. That is, there exists a function $\nabla g_i : \mathbb{R}^n \to \mathbb{R}^n$, the weak derivative of $g_i$, such that
    $$
g_i(y + z) - g_i(y) = \int_{0}^{1} z \cdot \nabla g_i(y + tz)\mathrm{d}t,
$$
for almost all $y, z \in \mathbb{R}^n$. Denote the Jacobian matrix of $g$ as $\nabla g \in \mathbb{R}^{n \times n}$, where the $i$-th row is equal to $\nabla g_i$.
In addition, assume that there exists $\alpha_0>0$ such that, for
$Z_{\alpha_0}\sim\calN(\mu,(1+\alpha_0)\sigma^2I_n)$,
$\mathbb{E}\norm{g(Z_{\alpha_0})}_2^4 < \infty$,
$\mathbb{E}\norm{\nabla g(Z_{\alpha_0})}_{\mathrm{F}}^2 < \infty$.

\end{assumption}

\begin{theorem}[Necessity and sufficiency of antithetic randomization]
\label{thm:alpha to 0}
Suppose that $(\omega^1,\ldots,\omega^K)$ satisfy Assumption~\ref{assump:normal-randomization} for some fixed $\rho$, and that $g$ satisfies Assumption~\ref{assump:weakdiff} and \sloppy{$\EE{\|g(Y)-Y\|^2_2}>0$}.
Then, as $\alpha\downarrow 0$,
\[
\mathbb{E} \left[\Var \!\left(\cv_{\alpha} \mid Y\right)\right] =\begin{cases}
O(1) &\text{ if } \rho=-\frac{1}{K-1},\\
\Theta(\frac{1}{\alpha}) &\text{ if } \rho>-\frac{1}{K-1}.
\end{cases}
\]
\end{theorem}

\begin{proposition}[A general construction of antithetic randomization schemes]
\label{lem: recipe:AR}
Let $Z\sim\calN(0,\sigma^2I_d)$ with $d\geq n$, and let $M=(M^1,\ldots,M^K)$ be an exchangeable family of random matrices in $\R^{n\times d}$, independent of $Z$, satisfying
\begin{enumerate}[label=(\roman*)]
    \item\label{cond:co-isometry} \emph{co-isometry:} $M^k(M^k)\tran=I_n$ for every $k\in[K]$;
    \item\label{cond:sum-zero-matrices} \emph{sum-zero:} $\displaystyle\sum_{k=1}^K M^k=0$.
\end{enumerate}
Define $\omega^k=M^kZ$ for $k\in[K]$.
Then the stacked vector $(\omega^1,\ldots,\omega^K)\in\R^{nK}$ satisfies Assumption~\ref{assump:normal-randomization} with $\rho=-1/(K-1)$ or, equivalently, the zero-sum constraint $\sum_{k=1}^K\omega^k=0$.
\end{proposition}

\begin{lemma}[Small-$\alpha$ limit]
\label{lem:smooth-small-noise-limit}
Suppose that $g$ satisfies Assumption~\ref{assump:weakdiff}, and let
$\omega^1,\ldots,\omega^K$ be independent of $Y$, with
$\omega^k\sim\calN(0,\sigma^2I_n)$ marginally for every $k$. Define
\begin{align}
R_\alpha
&:=
\frac1K\sum_{k=1}^K
\left\{
\|Y-g(Y+\sqrt\alpha\,\omega^k)\|_2^2
+
2\int_0^1
(\omega^k)\tran
\nabla g(Y+t\sqrt\alpha\,\omega^k)
\omega^k\,\rd t
\right\},\label{equ:def:R_alpha} \\ 
R_0
&:=
\|Y-g(Y)\|_2^2
+
\frac2K\sum_{k=1}^K
(\omega^k)\tran\nabla g(Y)\omega^k.\notag
\end{align}
Then
\[
R_\alpha\longrightarrow R_0
\qquad\text{in }L_2
\quad\text{as }\alpha\downarrow0.
\]
Consequently,
\[
\operatorname{Var}(R_\alpha\mid Y)
\longrightarrow
\operatorname{Var}(R_0\mid Y)
\qquad\text{in }L_1.
\]
\end{lemma}

\begin{lemma}[Variance decomposition]
\label{lem:reducible:variance}
Consider the randomization variables $\omega^k=M^kZ$ satisfying
\Cref{lem: recipe:AR} and independent of $Y$. Define
\[
H_g(Y)
:=
\frac12\left\{
\nabla g(Y)+\nabla g(Y)\tran
\right\},\qquad
D
:=
M^1(M^2)\tran+\frac{1}{K-1}I_n.
\]
Then
\begin{align*}
\lim_{\alpha\downarrow0}
\mathbb E\left[
\operatorname{Var}(\cv_\alpha\mid Y)
\right]
&=
\frac{8\sigma^4}{K-1}
\mathbb E\|H_g(Y)\|_{\mathrm F}^2
+
\frac{8(K-1)\sigma^4}{K}
\mathbb E\left[
\operatorname{tr}
\left\{
H_g(Y)D H_g(Y)D\tran
\right\}
\right].
\end{align*}
\end{lemma}

\begin{proof}[Proof of \Cref{lem:reducible:variance}]
By the sum-zero condition $\frac1K\sum_{k=1}^K\omega^k=0$,
the term $T_\alpha$ in the proof of
\Cref{thm:alpha to 0} vanishes and $\cv_\alpha=R_\alpha$, where $R_\alpha$ 
is defined in \Cref{equ:def:R_alpha}. By
\Cref{lem:smooth-small-noise-limit},
\[
\operatorname{Var}(\cv_\alpha\mid Y)
\longrightarrow
\operatorname{Var}\Bigl(
\frac2K\sum_{k=1}^K
(\omega^k)\tran\nabla g(Y)\omega^k
\, \mid \,Y
\Bigr)
\]
in $L_1$. Consequently,
\begin{align}
\lim_{\alpha\downarrow0}
\mathbb E\!\left[
\operatorname{Var}(\cv_\alpha\mid Y)
\right]
=
\mathbb E\Bigl[
\operatorname{Var}\Bigl(
\frac2K\sum_{k=1}^K
(\omega^k)\tran H_g(Y)\omega^k
\,\mid \,Y
\Bigr)
\Bigr],
\label{equ:rvar:limit}
\end{align}
where we used $x^\top A x=x^\top(A+A^\top)x/2$.

Condition on $Y$ and write $H=H_g(Y)$. Thus $H$ is a fixed symmetric
matrix in the following calculation. Set
\[
S_k:=(M^k)\tran H M^k,
\qquad
S:=\sum_{k=1}^KS_k.
\]
Since $\omega^k=M^kZ$,
\[
\frac2K\sum_{k=1}^K(\omega^k)\tran H\omega^k
=
\frac2K Z\tran S Z.
\]
Conditional on $M$, the matrix $S$ is symmetric and
$Z\sim\calN(0,\sigma^2I_d)$. The law of total variance therefore gives
\begin{align*}
V(H)
&:=
\operatorname{Var}\left(
\frac2K\sum_{k=1}^K(\omega^k)\tran H\omega^k
\right)\\
&=
\frac4{K^2}
\left\{
\mathbb E\!\left[
\operatorname{Var}(Z\tran S Z\mid M)
\right]
+
\operatorname{Var}\left(
\mathbb E[Z\tran S Z\mid M]
\right)
\right\}.
\end{align*}
For a symmetric matrix $S$,
\[
\mathbb E[Z\tran S Z\mid M]
=
\sigma^2\operatorname{tr}(S),
\qquad
\operatorname{Var}(Z\tran S Z\mid M)
=
2\sigma^4\|S\|_{\mathrm F}^2.
\]
Moreover, the co-isometry condition
$M^k(M^k)\tran=I_n$ gives
\[
\operatorname{tr}(S)
=
\sum_{k=1}^K
\operatorname{tr}\left\{
H M^k(M^k)\tran
\right\}
=
K\operatorname{tr}(H).
\]
Thus $\operatorname{tr}(S)$ is deterministic and
\begin{equation}
V(H)
=
\frac{8\sigma^4}{K^2}
\mathbb E\|S\|_{\mathrm F}^2.
\label{eq:V-corrected}
\end{equation}

Let $P_{jk}:=M^j(M^k)\tran$.
Since $H$ is symmetric,
\begin{align*}
\|S\|_{\mathrm F}^2
&=
\sum_{j,k=1}^K\operatorname{tr}(S_jS_k)=
\sum_{j,k=1}^K
\operatorname{tr}
\left(
H P_{jk}H P_{jk}\tran
\right).
\end{align*}
Using $P_{jj}=I_n$ and exchangeability,
\begin{equation}
\mathbb E\|S\|_{\mathrm F}^2
=
K\|H\|_{\mathrm F}^2
+
K(K-1)
\mathbb E\operatorname{tr}
\left(
H P_{12}H P_{12}\tran
\right).
\label{eq:S2-corrected}
\end{equation}

The sum-zero condition implies
\[
\sum_{k\neq1}P_{1k}
=
M^1\left(\sum_{k\neq1}M^k\right)\tran
=
-M^1(M^1)\tran
=
-I_n.
\]
By exchangeability,
$\mathbb E[P_{12}]=-\frac1{K-1}I_n$.
Writing
\[
P_{12}
=
-\frac1{K-1}I_n+D,
\qquad
\mathbb E[D]=0,
\]
we obtain
\begin{align*}
\mathbb E\operatorname{tr}
\left(
H P_{12}H P_{12}\tran
\right)
&=
\frac1{(K-1)^2}\operatorname{tr}(H^2)
+
\mathbb E\operatorname{tr}
\left(
H D H D\tran
\right),
\end{align*}
because the two terms linear in $D$ have expectation zero.

Substituting into \eqref{eq:S2-corrected} gives
\[
\mathbb E\|S\|_{\mathrm F}^2
=
\frac{K^2}{K-1}\|H\|_{\mathrm F}^2
+
K(K-1)
\mathbb E\operatorname{tr}
\left(
H D H D\tran
\right).
\]
Combining this identity with \eqref{eq:V-corrected},
\begin{align*}
V(H)
&=
\frac{8\sigma^4}{K-1}\|H\|_{\mathrm F}^2+
\frac{8(K-1)\sigma^4}{K}
\mathbb E\operatorname{tr}
\left(
H D H D\tran
\right).
\end{align*}
Finally, substitute $H=H_g(Y)$ and take expectation over $Y$ in
\eqref{equ:rvar:limit}.
\end{proof}

\end{document}
